Estado del arte / Marco teórico
Estado del arte
Existen enfoques de detección de “deepfakes” lo cuales tienen como objetivo principal la clasificación global binaria, creando una etiqueta según sus clases. El trabajo de Yagual2024 evaluó 3 tipos diferentes de arquitecturas de redes convolucionales dedicadas a la visión de computadora, las cuales fueron MobileNetV2, EfficientNetB0 y ResNet50, el desempeño de estas redes fue evaluado bajo un umbral de probabilidad dedicado a la clasificación binaria (Area Under Curve). Siguiendo con los trabajos de clasificación se encuentra Tiwari2025, utilizando una metodología hibrida combinando la arquitectura de ResNet-50 para la extracción de características robustas y posteriormente la clasificación ajustada de VGG-16 para la detección de deepfakes. Si bien ambos trabajos cumplen con sus objetivos comparten la deficiencia de no localizar espacialmente la alteración de la imagen, lo cual imposibilitan su uso en escenarios donde se requiere evaluar inconsistencias más detalladas.

Para poder resolver la generalidad que conlleva los detectores globales, investigaciones recientes han empezado a migrar hacia la perspectiva de localización sobre las alteraciones. En este sentido Deng2025 menciona como los clasificadores son incapaces de demostrar la relevancia de los componentes faciales manipulados, abordando un esquema de clasificación etiquetas múltiples por región afectada. En paralelo, Mashhadani2026 demuestra que la alteración de imágenes deja residuos visuales que debido al nivel de integración requieren ser ubicadas para tener una validación forense.

A diferencia de los trabajos anteriores, que resuelven la situación en donde la proporción de pixeles falsos contra píxeles reales es desbalanceada, con estrategias matemáticas como la función de perdida asimétrica polinomial (APL). En contraste en este proyecto el sistema busca centrarse en una segmentación semántica por medio de una red neuronal con una arquitectura basada en U-Net. Mediante su doble decodificador el modelo genera mascaras donde agrupa y delimita una región manipulada contra la real. Ademas, reduce el desbalanceo de clases usando una función de perdida compuesta de Entropía Cruzada Binaria (BCE), junto con métricas de segmentación semánticas (IoU y Dice) que se suman a la función de pérdida total.

Marco teórico

Manipulación facial y deepfakes

Definición operativa y tipología de manipulación facial

Mirsky2021 menciona que los deepfakes son contenido multimedia sintético o manipulado mediante técnicas de inteligencia artificial, donde la similitud visual se logra por modelos generativos capaces de reproducir rasgos faciales, texturas y patrones de iluminación con alta fidelidad y que pueden integrarse a flujos de edición que ocultan rastros evidentes mediante posprocesamiento, por ejemplo, compresión y el cambio de escala.

En particular, la manipulación facial puede organizarse como un conjunto de transformaciones sobre el rostro que abarcan:

 Reemplazo de identidad.
 Recreación o transferencia de expresiones.
 Ediciones localizadas del rostro sin sustituir por completo la identidad.
 Síntesis de un rostro completo sin una referencia. 


Técnicas como el reemplazo y la recreación tienden a introducir discontinuidades en contornos y coherencias internas del rostro, mientras que la síntesis puede producir coherencias globales altas, pero inconsistencias de textura en regiones específicas. MurguiaSerrano2025


La manipulación de rostros suele incluir una localización y emparejamiento facial, obtención de características, generación mediante modelos como redes generativas antagónicas (GAN) o modelos de difusión, y finalmente una fase de entremezclado y posprocesamiento para integrar el rostro manipulado con el contexto de la imagen real (color, grano, compresión y bordes). Estas etapas no solo producen una "imagen falsa", sino que también pueden dejar rastros residuales que se pueden ubicar localmente (texturas, micro bordes, patrones de ruido o inconsistencias de iluminación).

Detección de deepfakes: enfoques y taxonomía

Enfoques espacial, de frecuencia y aprendizaje profundo end-to-end


La detección de deepfakes se ha consolidado como un problema central del análisis forense digital debido a la rápida evolución de los generadores y a la creciente dificultad para distinguir contenido real de contenido sintético bajo condiciones de despliegue reales, con especial énfasis en la generalización y en la degradación introducida por compresión y reescalado.

Debido al acelerado desarrollo de generadores, la deteccion se puede estructurar en funcion del tipo de evidencia que estas pueden modelar, por ello se adopta una taxonomia de cuatro enfoques de alto nivel:

 Enfoque espacial.
 Enfoque en frecuencia.
 Enfoque basado en aprendizaje profundo (end-to-end).
 Enfoques híbridos. 

Esta organización no pretende ser excluyente, ya que diversos modelos contemporáneos combinan señales y representaciones; sin embargo, permite sintetizar la intuición forense detrás de cada familia y, en consecuencia, discutir sus fallas típicas ante el cambio de dominio (conjunto de datos, generador y postprocesado) MurguiaSerrano2025, Banerjee2025.

En el enfoque de frecuencia, la hipótesis es que los procesos de generación introducen artefactos espectrales que no son necesariamente evidentes en el espacio de píxel, pero sí se manifiestan como anomalías en bandas de alta frecuencia o como distribuciones espectrales atípicas. Parte de la evidencia reportada en la bibliografía vincula estas huellas con operaciones internas de los generadores, por ejemplo, etapas de sobremuestreo en arquitecturas GAN, que dejan patrones característicos en el espectro. En consecuencia, se utilizan transformadas (p. ej., Fourier, DCT o wavelets) y descriptores espectrales para separar contenido real de contenido sintético. Sin embargo, delimitante importante es que las huellas pueden variar entre generadores y, por lo tanto, indcuir cierta sensibilidad a patrones específicos del modelo de síntesis o del dominio de entrenamiento, reduciendo la capacidad frente a escenarios mas realista que suelen ser mas variados. Tan2024


En los enfoques de basados de aprendizaje profundo y en el dominio espacial, la detección se centra en inconsistencias observables en el dominio de píxel, es decir, irregularidades locales o globales que emergen de la unión, la edición y la síntesis facial. Habitualmente se buscan discontinuidades en contornos, incoherencias en iluminación, texturas atípicas, patrones de ruido residuales o discrepancias entre la región facial y su contexto. En implementaciones modernas, estas señales se capturan mediante extractores aprendidos, que identifican patrones discriminativos asociados a manipulaciones y, en variantes orientadas a localización, habilitan mapas de activación o máscaras que aproximan la región alterada, las arquitecturas CNN y variantes con mecanismos de atención o módulos complementarios, entrenadas en clasificación (real/falso) o en formulaciones que incorporan supervisión espacial cuando se dispone de anotaciones. Su ventaja práctica es la capacidad de modelar patrones complejos y no lineales. No obstante, No obstante, al depender de huellas que pueden ser atenuadas por posprocesamiento, el desempeño suele degradarse en escenarios con compresión, baja resolución u oclusiones, además de mostrar sensibilidad a sesgos de captura y a variaciones de escena. Banerjee2025, Alrashoud2025


Finalmente, los enfoques híbridos surgen como respuesta directa a la fragilidad de un único tipo de evidencia. De manera general, un enfoque híbrido integra representaciones espaciales y espectrales, o combina señales forenses (p. ej., inconsistencias de compresión o señales residuales) con características aprendidas por redes profundas, mediante arquitecturas multirrama y mecanismos de fusión. La intuición es que la manipulación deja rastros en múltiples niveles: algunos son locales en el píxel, otros se expresan como firmas en frecuencia y otros se manifiestan como discrepancias de coherencia global; por lo tanto, la combinación explícita puede aumentar sensibilidad y robustez ante variaciones. En la bibliografía reciente se observa un énfasis en modelar interacciones espacio–frecuencia y en diseñar estrategias orientadas a mejorar la generalización, precisamente porque los artefactos aislados tienden a perder estabilidad frente a generadores nuevos y frente a interferencia por postprocesado. Li2025, Qiao2025
 
En este marco, la generalización se entiende como la capacidad del detector para mantener desempeño ante cambios de dominio: nuevos generadores, nuevas distribuciones de datos y nuevas condiciones de captura y publicación. La evidencia reciente subraya que, incluso cuando los modelos obtienen resultados altos en conjuntos académicos, su desempeño puede caer de forma pronunciada en datos "in-the-wild", recolectados de plataformas y escenarios contemporáneos. Adicionalmente, la compresión, el reescalado y otros posprocesos actúan como perturbaciones que suprimen o distorsionan huellas forenses, afectando tanto señales espaciales (pérdida de microtextura) como señales espectrales (alteración de distribuciones de frecuencia). Hu2024, Li2025, Chandra2025

Clasificación global vs localización

Limitaciones de la clasificación global ante manipulaciones locales

La detección de deepfakes se podría conceptualizar como un problema de clasificación global, en donde un modelo asigna una etiqueta binaria (real/falso) a la imagen completa. En escenarios reales, las alteraciones llegan a modificar únicamente regiones específicas del rostro (por ejemplo, ojos, boca o contornos faciales). Cuando la alteración está delimitada a una región, los modelos de clasificación global pueden no detectar señales discriminativas, especialmente después de procesos de compresión, reescalado o suavizado que reducen las huellas forenses visibles. Esto aumenta la probabilidad de falsos negativos, ya que la imagen conserva gran parte de sus píxeles sin manipular.Rossler_2019_ICCV


Reformulación como tarea de localización


El problema abordado en este trabajo se reformula como una tarea de localización de manipulación, donde el objetivo no es únicamente determinar si una imagen es falsa, sino identificar qué regiones específicas han sido alteradas. En lugar de producir una etiqueta global, el modelo genera un mapa espacial o máscara que asigna a cada píxel una probabilidad de pertenecer a una zona manipulada. Esta metodología integra a la segmentación semántica para poder generar un mapa espacial ayudando a identificar dónde se realizó una alteración.

La generación de una máscara proporciona evidencia visual interpretable, permite observar la extensión de la manipulación, para posteriormente facilitar los análisis centrados exclusivamente en las regiones sospechosas. Bappy_2017_ICCV

Fundamentos de CNN para visión por computadora


Redes neuronales convolucionales


goodfellow2016deep describen las redes neuronales convolucionales como un tipo de red neuronal specializada en el procesamiento de datos que tiene una topología conocida similar a una cuadrícula. Mencionan particularmente a las imágenes, las cuales pueden considerarse como una cuadrícula bidimensional de píxeles. El nombre "red neuronal convolucional" indica que la red emplea una operación matemática llamada convolución. La convolución es un tipo specializado de operación lineal. Por su parte las redes neuronales convolucionales son simplemente redes neuronales que utilizan la convolución en lugar de la multiplicación matricial general en al menos una de sus capas. 


La operación de convolución


La operación de convolución se define como una operación matemática entre dos funciones que produce una tercera función que representa la cantidad de superposición entre ellas. En términos continuos puede expresarse como:
 

Donde:

 	x representa la señal de entrada o los datos originales. En el contexto de las redes neuronales convolucionales, esta señal corresponde a los valores de los píxeles de una imagen
 	w representa el filtro o kernel de convolución en el contexto de CNN el cual es una matriz de valores que se utilizan para detectar características específicas en la imagen (Bordes, texturas, patrones, etc.)
 	* representa la operación de convolución entre la señal de entrada x y el filtro w
 	t representa la posición o desplazamiento en la señal o imagen donde se está aplicando el filtro
 	s(t) es el resultado de la convolución, es decir, una señal nueva resultante del proceso de convolución o un mapa de activación en el contexto de las CNN obtenido después de aplicar un filtro a una señal de entrada.



En el contexto de redes neuronales convolucionales, esta operación se utiliza para aplicar filtros sobre imágenes con el objetivo de detectar patrones locales como bordes, texturas o formas. Cada filtro se desplaza sobre la imagen generando mapas de activación que resaltan dichas características.

Sin embargo, en aplicaciones de aprendizaje automático y visión por computadora, los datos utilizados no son continuos sino discretos. Las imágenes digitales, por ejemplo, se representan como matrices de valores numéricos correspondientes a píxeles. Debido a esto, la operación de convolución se expresa mediante una forma discreta, en la cual la integral es reemplazada por una suma.


Debido a que las imágenes digitales se pueden representar como matrices bidimensionales de píxeles, para llevar a cabo el proceso de convolución se aplica la misma fórmula, pero adaptada a dos dimensiones:


Aunque el término convolución es ampliamente utilizado en redes neuronales convolucionales, en la práctica la operación implementada en la mayoría de los frameworks de aprendizaje profundo corresponde técnicamente a una correlación cruzada (cross-correlation). La principal diferencia entre ambas operaciones radica en que la convolución implica una inversión del filtro antes de realizar el desplazamiento sobre la señal de entrada, mientras que en la correlación el filtro se aplica directamente sin dicha inversión. A pesar de esta diferencia matemática, en el contexto del aprendizaje profundo ambos términos suelen utilizarse de manera intercambiable, ya que los filtros son aprendidos automáticamente durante el entrenamiento de la red.


Filtros, mapas de activación y jerarquía de características.


En las redes neuronales convolucionales, los filtros o kernels son pequeñas matrices de valores que se utilizan para detectar características específicas dentro de una imagen. Estos filtros se desplazan sobre la matriz de entrada aplicando la operación de convolución o correlación en cada posición. Durante el proceso de entrenamiento, los valores de los filtros se ajustan automáticamente mediante algoritmos de optimización, permitiendo que la red aprenda a identificar patrones relevantes en los datos.


 


El resultado de aplicar un filtro sobre una imagen de entrada se denomina mapa de activación (feature map). Este mapa representa la respuesta del filtro en cada posición de la imagen, resaltando aquellas regiones donde el patrón detectado por el filtro está presente. Cada filtro genera un mapa de activación diferente, lo que permite a la red neuronal capturar múltiples características visuales dentro de la imagen.


Las redes neuronales convolucionales aprenden una jerarquía de características visuales a través de sus diferentes capas. En las primeras capas, los filtros suelen detectar características simples como bordes y orientaciones. En capas intermedias, la red comienza a identificar patrones más complejos como texturas o formas. Finalmente, en las capas más profundas se reconocen estructuras más complejas, como partes de objetos o configuraciones faciales completas. Esta capacidad de aprender representaciones jerárquicas permite a las CNN analizar imágenes de forma eficiente en tareas de reconocimiento y clasificación.
Esta capacidad resulta especialmente útil en la detección de deepfakes, ya que permite identificar inconsistencias visuales sutiles en imágenes manipuladas, tales como anomalías en texturas faciales, iluminación o patrones de compresión.


En las redes neuronales convolucionales, esta operación se aplica repetidamente a través de múltiples capas convolucionales. Cada capa utiliza varios filtros que son aprendidos automáticamente durante el proceso de entrenamiento, permitiendo a la red detectar diferentes tipos de características visuales. Como resultado, se generan múltiples mapas de activación que capturan patrones relevantes dentro de la imagen de entrada.


Pooling y stride.

En las redes neuronales convolucionales, el pooling es una operación mediante la cual se busca reducir la dimensionalidad o tamaño de los mapas de activación generados por las capas convolucionales. Esta reducción permite disminuir el número de parámetros que utiliza el modelo y el costo computacional, al mismo tiempo que se conserva la información más relevante para las características que anteriormente se detectaron. Este proceso consiste en dividir el mapa de activación en pequeñas regiones y aplicar una función de agregación sobre cada una de ellas dando como resultado un mapa de activación más pequeño que conserva las características más importantes detectadas por los filtros de la red.


El max pooling es uno de los métodos de pooling más utilizados en redes neuronales convolucionales. En este método, para cada región del mapa de activación se selecciona únicamente el valor máximo. De esta manera, se conserva la característica más significativa detectada por el filtro en esa región específica.
Este mecanismo permite que la red neuronal mantenga las activaciones más relevantes y reduzca la información irrelevante, lo que ayuda a mejorar la eficiencia del modelo.


Otro método utilizado es el average pooling, el cual calcula el valor promedio dentro de cada región del mapa de activación. Aunque este método también reduce la dimensionalidad de los datos, en la práctica el max pooling suele utilizarse con mayor frecuencia en arquitecturas modernas de redes neuronales convolucionales debido a su mejor desempeño en muchas tareas de visión por computadora.


En la operación de convolución, el stride (comúnmente traducido como paso o zancada) es un hiperparámetro estructural que determina la cantidad de unidades espaciales o píxeles que el filtro (kernel) avanza sobre la matriz de entrada en cada iteración.
En las operaciones de convolución estándar, se asume un desplazamiento unitario (stride = 1), lo que significa que el filtro se mueve un píxel a la vez de izquierda a derecha y de arriba hacia abajo, procesando la imagen de manera altamente superpuesta.
 Sin embargo, al incrementar este valor (por ejemplo, un stride de 2), el filtro "salta" posiciones, desplazándose de dos en dos píxeles.
El impacto principal del stride se ve reflejado en las dimensiones espaciales del mapa de activación resultante una vez terminado el proceso de convolución, si el stride es grande, las dimensiones serán menores, actuando como un mecanismo de submuestreo y se denota de la siguiente manera:

Funciones de activación.


En las redes neuronales convolucionales, después de realizar la operación de convolución, se aplica una función de activación. Estas funciones introducen no linealidad en el modelo, lo cual permite que la red neuronal pueda aprender representaciones más complejas de los datos. Sin la incorporación de funciones no lineales, una red compuesta únicamente por operaciones lineales sería equivalente a una sola transformación lineal, limitando significativamente su capacidad de aprendizaje.
Las funciones de activación se aplican a los valores generados por las capas convolucionales, produciendo mapas de activación que posteriormente son procesados por las siguientes capas de la red. De acuerdo con goodfellow2016deep, las funciones de activación permiten que las redes neuronales profundas modelen relaciones complejas presentes en los datos y facilitan el aprendizaje de representaciones jerárquicas en tareas de visión por computadora.

Una de las funciones de activación más utilizadas en redes neuronales convolucionales es la Rectified Linear Unit (ReLU). Esta función se caracteriza por su simplicidad y eficiencia computacional, ya que transforma los valores negativos en cero y mantiene los valores positivos sin modificación.

Según goodfellow2016deep, el uso de la función ReLU facilita el entrenamiento de redes neuronales profundas, ya que ayuda a reducir el problema del gradiente desvanecido y permite una convergencia más rápida durante el proceso de aprendizaje.

Aun con la eficiencia de ReLU en las capas ocultas para la extracción de características, la salid de la red ocupa un enfoque diferente, pues las clasificación binaria como la generación de mascaras en problemas de segmentación semántica mediante arquitecturas como U-Net Ronneberger2015, por ellos la salid del modela se interpreta como una probabilidad, la función de activación sigmoide, comprime el valor real de la de entrada en un intervalo continuo entre el 0 y 1 lo cual según goodfellow2016deep resulta determinante para calcular la probabilidad de pertenecer a una clase.

Segmentación semántica y formulación del problema.

Artefactos e inconsistencias en deepfake.

El uso desmedido de redes generativas (GAN) para producir deepfakes ha conducido a una proliferación de imágenes y con ello, al aumento del realismo. Sin embargo, aún con todos estos avances se pueden detectar por las técnicas correctas artefactos visuales que delatan su intervención. “FaceForensics++: Learning to Detect Manipulated Facial Images de Andreas Rossler” demostró que aún existe irregularidades en textura de piel asi como iluminación y sombras.

Dentro del dominio contextual de un rostro, la mayor inconsistencia presente de texturas se presenta cuando se realiza un suavizado en la piel o bien cuando estén perdida de imperfecciones que la piel debería de conservar y más con los avances de calidad digital al momento de capturar contenido. Este tipo de efectos surgen cuando en la generación de la manipulación se le da prioridad a una coherencia global sobre las particulares detalladas (Ejemplo: Un tono de piel informe por sobre la generación de líneas de expresión) lo que provoca transiciones notorias entre el rostro alterado y la imagen real, las zonas que más se delatan son el cabello y el contorno facial. Otro tipo de incoherencia visual que se presenta con regularidad son la dualidad de sombra y luz que resaltan por no coincidir con la imagen real.

Altas frecuencias y residuos de textura.

Un aspecto para tomar en cuenta es el análisis de alta frecuencia, el cual suele abarcar información acerca de bordes, detalles fino o micro texturas. Los cuales se alteran durante la generación de un “deepfake”. Detección como el propuesto como “MesoNet: a Compact Facial Video Forgery Detection Network de Darius Afchar” mencionan que se puede tomar ventaja de los residuos de textura para extraer características o rastros de manipulación. En varios casos la region manipulada presenta un conjunto de frecuencias distintas a las imágenes reales, permitiendo detección que serían difícil para la percepción del ojo del ser humano 

Segmentación semántica.

Las técnicas de segmentación semántica permiten asignar una etiqueta a cada píxel de una imagen según su contexto. A diferencia de un método tradicional de clasificación que genera una etiqueta global, la segmentación permite determinar ubicación dentro las imágenes, dicho enfoque resulta útil en un análisis forense permitiendo la localización espacial.

Formulación del problema como segmentación binaria.

El objetivo del modelo es aprender una función que a partir de una imagen del rostro pueda generar un mascara binaria que indique regiones manipuladas, De esta forma los pixeles pertenecientes a una clase se pueden representar asignándole el calor 0, caso contrario los pixeles con de regiones de otras clases se podrá asignar el valor 1. Dicho formulación permitirá localizar con precisión las áreas de la alteración.


Arquitectura U-Net

Qué es U-Net

La red neuronal U-Net es una red convolucional diseñada para tareas de segmentación, propuesta por Ronneberger2015 en 2015 fue diseñada para usarse en contextos de medicina donde se requirió precisión en imágenes de tipo biomédicas ya que se requería localizar estructuras específicas píxel a píxel. Este modelo se diferencia al resto de clasificadores globales por tener una predicción mas especifica, permitiendo localizar regiones donde la tarea exige una ubicación del objeto a detectar, llámese alteración u objeto. Su nombre se deriva de la característica de su arquitectura ya que en la etapa de contracción o codificación seguida de por la etapa de expansión o decodificación genera un peculiar y reconocible “U”, su arquitectura permite la amplia detección desde características generales hasta los detalles más finos.

Componentes principales de la arquitectura

La U-Net se compone de 3 elementos particulares, el de codificador, decodificador y los saltos de conexión. 
El primero se encarga de extraer las características de la imagen de entrada, esta diseñada con una serie de capas convolucionales acompañadas con operaciones de max pooling que reducen de manera progresiva la resolución mientras aumentan la profundidad de los mapas de características extraídos. Su objetivo principal es juntar todo el contexto e información semántica, logrando que el modelo identifique patrones complejos. Cada bloque que constituye el codificador está compuesto generalmente por:

 Dos convoluciones de 3x3
 Una función de activación
 Una operación de max pooling


La reducción permite el aprendizaje por abstracción.
El Decodificador tiene la función de volver a construir la resolución original de la imagen para generar el mapa de segmentación, es dicha etapa se utilizan operaciones de up-smapling o convolución transpuesta, aumentando progresivamente la resolución. El bloque de un decodificador se constituye por:

   
 Up-sampling
 Concatenación de mapas de características
 Dos convoluciones de 3x3


El decodificador recuera los detalles del espacio que se pueden llegar a perder del proceso anterior.

Los saltos de conexión son de los elementos mas importantes de la arquitectura U-Net, las conexiones se enlazan entres las capas del codificador y el decodificador, cada una unida a su capa correspondiente, concatenando los mapas de características con el propósito de preservar la mayor cantidad de información espacial de alta resolución que de no hacer estas conexiones se perderían durante el proceso de reducción del codificador, gracias a estas conexiones el modelo puede combinar información global hasta los detalles finas.


Diferencia entre clasificación y segmentación
El objetivo de un modelo de clasificación es asignar una etiqueta general a la imagen, para llevar a cabo la tarea, las redes sueles reducir la resolución de manera progresiva, sin embargo cuando el objetivo cambia hacia la segmentación el objetivo del modelo es construir un mapa de predicción denso en el cual el pixel tiene la etiqueta, implicando que las características globales se deban de mantener o reconstruirse para producir la información espacial y predicción en cada pixel.


